\scanpage{155}
% archive foliation: 154.
Normalité des extensions quadratiques
\[
A[S]/(S^2-a)=B
\]
On suppose d'abord $A$ \unsure{anneau jauge}, soit $p$ sa car. résiduelle.
$p\ne2$, alors $B$ normal ssi $a\notin\mathfrak m^2$ --- et cette condition
est nécessaire en tous cas (y compris $p=2$). Si $p=2$, distinguons
deux cas:

\margin{\uline{NB} si $p\ne2$, alors $B$ est local ssi on a
$(a\unsure{\in|\notin}\mathfrak m,\text{ et})$ \uline{ou}
\deleted{\unsure{alors}} $\bar a\notin k^2$}

a) $a\in\mathfrak m$ alors $B$ normal ssi $a\notin\mathfrak m^2$

b) $a\notin\mathfrak m$ i.e. $a$ une unité. Soit $\bar a$ son image dans
$k:=k(A)=A/\mathfrak m$.

$1^\circ)$ $\bar a\notin k^2$, alors $B$ est normal (est un \unsure{anneau jauge} de
corps résiduel $k[\sqrt{\bar a}]$).

$2^\circ)$ $\bar a\in k^2$, soit $\bar a=\bar\alpha^2$, où $\alpha\in A$ déterminé mod
$\mathfrak m\deleted{\unsure{B}}$ --- donc
$a=\alpha^2\deleted{\unsure{+}}-b$, avec $b\inserted{={}\alpha^2-a}\in\mathfrak m$.
Posons $T=S-\alpha$, alors
\[
S^2-a=(T+\alpha)^2-a=T^2+2\alpha T+b
\]
\[
B\simeq A[T]/(T^2+2\alpha T+b),\quad
\widehat B\simeq\underbrace{A[[T]]}_{C}/(T^2+2\alpha T+b)
\]
\margin{\uline{NB} $\widetilde b\deleted{\unsure{\mapsto}}\in\mathfrak m/\mathfrak m^2$
ne dépend que de $a$, et pas du choix de $\alpha$}
\margin{NB $2\in\mathfrak m\subset\mathfrak r(C)$,\quad $T\in\mathfrak r(C)$}
$T^2,2\alpha T$ sont dans $\unsure{\mathfrak r|\mathfrak m}(C)^2$, donc
$T^2+2\alpha T+b\notin(\mathfrak r(C))^2$
ssi $b\notin\mathfrak r(C)^2$ i.e $b\notin\mathfrak m^2$, i.e.
$\widetilde b\ne0$ dans $\mathfrak m/\mathfrak m^2$.

Donc les cas de normalité sont les suivants
\[
a\in\mathfrak m\setminus\mathfrak m^2,\quad
\text{\uline{ou} $a$ de la forme $\alpha^2-b$ avec $\alpha$ unité, $b\in\mathfrak m\setminus\mathfrak m^2$}
\]
\[
\text{\uline{ou}}\quad a\notin\mathfrak m,\quad\bar a\notin k^2
\]
\margin{Dans les cas de $B$ \uline{non} normal, on a $a\in\mathfrak m^2$, \uline{ou}
$a\bmod\mathfrak m^2$ le carré d'une unité de $A/\mathfrak m^2$.
\uline{Ce dernier cas correspond à une \unsure{singularité} quadratique ordinaire}.}

\uline{Application} aux rev. quadratiques de $\mathbb P^1_{\mathbb Z}$
\inserted{(ou de $\mathbb E^1_{\mathbb Z}$)}
génériquement \unsure{non étales}
dans chaque fibre au dessous de $\operatorname{Spec}\mathbb Z$.

Le revêtement \deleted{\unsure{normal}} \deleted{\unsure{en dessus de $\eta$}}
\inserted{quadratique} du pt gén $\xi$ correspond à une ext.
quadratique de $\mathbb Q(t)$, donc décrit par un élt de
$\mathbb Q(t)^*/\mathbb Q(t)^{*2}$.
Or $\unsure{\mathbb Q(t)^*/\{\pm1\}}$ a une base formée des éléments suivants

a) Les nbs premiers $p\in\mathbb P$\quad
$(\mathbb P=\{2,3,5,7\ldots\})$

b) Les polynomes $P=a_0t^n+a_1t^{n-1}+\cdots+a_n\in\mathbb Z[t]$, avec:
$a_0>0$, \deleted{\unsure{et}} les $a_i$ premiers entre eux
(i.e. $\forall p\in\mathbb P$, $P\bmod p\in\mathbb Z_p[t]$ est $\ne0$),
$P$ premier dans $\mathbb Q[T]$.

Ainsi un élément de $\mathbb Q(t)^*/\mathbb Q(t)^{*2}$ est représenté de façon
unique par un produit
\[
P=\underbrace{\pm}_{\varepsilon}
\underbrace{p_1p_2\cdots p_n}_{q}
\underbrace{P_1\cdots P_\ell}_{Q}
\tag{1}
\]
où les $p_i\in\mathbb P$ sont distincts, et les $P_j$ comme dans b) aussi.
Cela signifie aussi qu'on a
\[
P=\pm qQ\qquad
\left\{
\begin{aligned}
&q\in\mathbb N\text{ ``quadratfrei''}\\
&Q\in\mathbb Z[t]\text{ à coeff. dominant }>0\\
&\text{``quadratfrei'' dans $\mathbb Q[t]$}\\
&\text{à coeff. premiers entre eux}
\end{aligned}
\right.
\]
